% GENERATED FILE. Edit canonical lessons, then run npm run generate:books.
% canonical-source-hash: fnv1a64:700246e76f28c597
% canonical-lessons: FR-C410-bouchon, FR-C410-embouteillage, FR-C410-deviation, FR-C410-aire, FR-C410-doubler

\chapter{Bouchons sur l'autoroute}
\label{ch:fr-bouchons-sur-l-autoroute}
\hlchaptermodality{410}

\begin{chapteropening}
\textbf{By the end of this chapter:} \emph{I can talk about heavy traffic: the jam, the gridlock, the diversion, the rest area and overtaking.}

Complete the last lesson of chapter 410: I can talk about heavy traffic: the jam, the gridlock, the diversion, the rest area and overtaking.
\end{chapteropening}

\section[le bouchon]{le bouchon — the traffic jam; the cork}

\label{lesson:FR-C410-bouchon}



\begin{quote}
Before the new one: say the French for the road surface, then the French for to turn aside.
\end{quote}

\subsection*{You'll want to know: le bouchon}
\textbf{le bouchon} — \textquotedblleft{}the traffic jam; the cork\textquotedblright{}.

\begin{quote}
Il y a dix kilomètres de bouchon avant le tunnel.
\end{quote}

\begin{quote}
\emph{There is a ten-kilometre jam before the tunnel.}
\end{quote}

\emph{Un bouchon} is also a cork for a bottle.

\begin{cousinweb}
\emph{bouchon} is built on Middle French \emph{bousche} 'handful of straw' (EtymDB): first a twist of straw that stopped up a bottle, now a cork and a road blocked by cars.
\end{cousinweb}

\begin{grammarlens}[title={What you've built}]
One more word for telling, explaining and describing at B1.
\end{grammarlens}

\subsection*{Your turn}
\begin{itemize}
\raggedright
  \item \emph{Say it:} \emph{le bouchon}
  \item \emph{Say it:} \emph{le bouchon}, once more
  \item \emph{Say it:} a sentence of your own with \emph{le bouchon}
  \item \emph{Recall:} say the French for the road surface, then the French for to turn aside, then say \emph{le bouchon} again
\end{itemize}

\begin{tcolorbox}[breakable,colback=teal!4,colframe=teal!35!black,title={Before you move on}]
What does le bouchon mean? (\textquotedblleft{}The traffic jam; the cork\textquotedblright{}.) Say it once more.
\end{tcolorbox}
\section[l'embouteillage]{l'embouteillage — the gridlock, the traffic jam}

\label{lesson:FR-C410-embouteillage}



\begin{quote}
Before the new one: say the French for to turn aside, then the French for the traffic jam.
\end{quote}

\subsection*{You'll want to know: l'embouteillage}
\textbf{l'embouteillage} — \textquotedblleft{}the gridlock, the traffic jam\textquotedblright{}.

\begin{quote}
Le vendredi soir, les embouteillages sont terribles.
\end{quote}

\begin{quote}
\emph{On Friday evenings the traffic jams are terrible.}
\end{quote}

A masculine noun; close to \emph{bouchon}.

\begin{cousinweb}
\emph{embouteillage} is built on \emph{embouteiller} 'to put in bottles', from \emph{bouteille} 'bottle', which goes back to Late Latin \emph{butticula}, a small \emph{buttis} 'cask' (EtymDB): cars stuck as in the neck of a bottle.
\end{cousinweb}

\begin{grammarlens}[title={What you've built}]
One more word for telling, explaining and describing at B1.
\end{grammarlens}

\subsection*{Your turn}
\begin{itemize}
\raggedright
  \item \emph{Say it:} \emph{l'embouteillage}
  \item \emph{Say it:} \emph{l'embouteillage}, once more
  \item \emph{Say it:} a sentence of your own with \emph{l'embouteillage}
  \item \emph{Recall:} say the French for to turn aside, then the French for the traffic jam, then say \emph{l'embouteillage} again
\end{itemize}

\begin{tcolorbox}[breakable,colback=teal!4,colframe=teal!35!black,title={Before you move on}]
What does l'embouteillage mean? (\textquotedblleft{}The gridlock, the traffic jam\textquotedblright{}.) Say it once more.
\end{tcolorbox}
\section[la déviation]{la déviation — the diversion, the detour (for traffic)}

\label{lesson:FR-C410-deviation}



\begin{quote}
Before the new one: say the French for the traffic jam, then the French for the gridlock.
\end{quote}

\subsection*{You'll want to know: la déviation}
\textbf{la déviation} — \textquotedblleft{}the diversion, the detour (for traffic)\textquotedblright{}.

\begin{quote}
À cause des travaux, il y a une déviation par le village.
\end{quote}

\begin{quote}
\emph{Because of the roadworks there is a diversion through the village.}
\end{quote}

Signs say \emph{Déviation} with an arrow.

\begin{cousinweb}
\emph{déviation} is a learned borrowing from Late Latin \emph{dēviātiō}, from \emph{dēviāre} 'to leave the road', built on \emph{via} 'way' (Lewis \& Short).
\end{cousinweb}

\begin{grammarlens}[title={What you've built}]
One more word for telling, explaining and describing at B1.
\end{grammarlens}

\subsection*{Your turn}
\begin{itemize}
\raggedright
  \item \emph{Say it:} \emph{la déviation}
  \item \emph{Say it:} \emph{la déviation}, once more
  \item \emph{Say it:} a sentence of your own with \emph{la déviation}
  \item \emph{Recall:} say the French for the traffic jam, then the French for the gridlock, then say \emph{la déviation} again
\end{itemize}

\begin{tcolorbox}[breakable,colback=teal!4,colframe=teal!35!black,title={Before you move on}]
What does la déviation mean? (\textquotedblleft{}The diversion, the detour (for traffic)\textquotedblright{}.) Say it once more.
\end{tcolorbox}
\section[l'aire]{l'aire — the rest area (on a motorway); the area}

\label{lesson:FR-C410-aire}



\begin{quote}
Before the new one: say the French for the gridlock, then the French for the diversion.
\end{quote}

\subsection*{You'll want to know: l'aire}
\textbf{l'aire} — \textquotedblleft{}the rest area (on a motorway); the area\textquotedblright{}.

\begin{quote}
Nous mangeons sur une aire de repos.
\end{quote}

\begin{quote}
\emph{We eat at a motorway rest area.}
\end{quote}

A feminine noun: \emph{une aire de repos}, \emph{une aire de jeux}.

\begin{cousinweb}
\emph{aire} is inherited from Latin \emph{ārea} 'open space, threshing floor' (EtymDB).
\end{cousinweb}

\begin{grammarlens}[title={What you've built}]
One more word for telling, explaining and describing at B1.
\end{grammarlens}

\subsection*{Your turn}
\begin{itemize}
\raggedright
  \item \emph{Say it:} \emph{l'aire}
  \item \emph{Say it:} \emph{l'aire}, once more
  \item \emph{Say it:} a sentence of your own with \emph{l'aire}
  \item \emph{Recall:} say the French for the gridlock, then the French for the diversion, then say \emph{l'aire} again
\end{itemize}

\begin{tcolorbox}[breakable,colback=teal!4,colframe=teal!35!black,title={Before you move on}]
What does l'aire mean? (\textquotedblleft{}The rest area (on a motorway); the area\textquotedblright{}.) Say it once more.
\end{tcolorbox}
\section[doubler]{doubler — to overtake; to double}

\label{lesson:FR-C410-doubler}



\begin{quote}
Before the new one: say the French for the diversion, then the French for the rest area.
\end{quote}

\subsection*{You'll want to know: doubler}
\textbf{doubler} — \textquotedblleft{}to overtake; to double\textquotedblright{}.

\begin{quote}
Il est dangereux de doubler dans un tunnel.
\end{quote}

\begin{quote}
\emph{It is dangerous to overtake in a tunnel.}
\end{quote}

\emph{Dépasser} also means to overtake.

\begin{cousinweb}
\emph{doubler} is inherited through Old French \emph{dobler} from Late Latin \emph{dūplāre} 'to double', from \emph{duplus} 'double' (EtymDB).
\end{cousinweb}

\begin{grammarlens}[title={What you've built}]
One more word for telling, explaining and describing at B1.
\end{grammarlens}

\subsection*{Your turn}
\begin{itemize}
\raggedright
  \item \emph{Say it:} \emph{doubler}
  \item \emph{Say it:} \emph{doubler}, once more
  \item \emph{Say it:} a sentence of your own with \emph{doubler}
  \item \emph{Recall:} say the French for the diversion, then the French for the rest area, then say \emph{doubler} again
\end{itemize}

\begin{tcolorbox}[breakable,colback=teal!4,colframe=teal!35!black,title={Before you move on}]
What does doubler mean? (\textquotedblleft{}To overtake; to double\textquotedblright{}.) Say it once more.
\end{tcolorbox}
